\section{Introduction}
\label{sec:introduction}

For a convex body $K\subset\R^n$, its projection body $\Pi K$ is the
origin-symmetric convex body whose support function is
\[
 h_{\Pi K}(u)=\operatorname{vol}_{n-1}
       \bigl(\operatorname{proj}_{u^\perp}K\bigr),
       \qquad u\in S^{n-1}.
\]
Here a convex body is compact, convex, and has nonempty interior, and
$h_L(x)=\max_{y\in L}\ip{x}{y}$ denotes its support function.
Thus $\Pi K$ records the volumes of the orthogonal shadows of $K$ in
all directions. Write $|L|=\operatorname{vol}_n(L)$ and let $B_2^m$
be the Euclidean unit ball in $\R^m$, with volume $\kappa_m$.
The normalized projection volume is
\[
 R_n(K)=\frac{|\Pi K|}{|K|^{n-1}}.
\]
It is invariant under translations and invertible linear maps: the
classical contravariance identity is
$\Pi(TK)=|\det T|T^{-t}\Pi K$; see
\cite[Equation~(1)]{Ludwig2002}. We verify this identity directly in
Section~\ref{sec:conclusion}.

Petty's projection-volume conjecture asserts, for $n\ge3$, that
$R_n$ is minimized precisely by ellipsoids \cite{Petty1971}.
The ball satisfies $\Pi B_2^n=\kappa_{n-1}B_2^n$, so the proposed
minimum is fixed by the Euclidean case. This is an affine isoperimetric
problem in which directional projection data are assembled into a
second body before its volume is measured. The classical Petty
inequality for the \emph{polar} of the projection body concerns a
different functional; see \cite[Section~18.6]{Gardner2002}.
We prove the unpolarized conjecture in every dimension at least four.

\begin{theorem}\label{thm:main}
For every integer $n\ge4$ and every convex body $K\subset\R^n$,
\begin{equation}\label{eq:main}
 \frac{|\Pi K|}{|K|^{n-1}}
 \ge \kappa_{n-1}^{\,n}\kappa_n^{\,2-n}.
\end{equation}
Equality holds if and only if $K=a+TB_2^n$ for some $a\in\R^n$ and some
invertible linear map $T$.
\end{theorem}

The theorem includes bodies without symmetry or boundary regularity.
Its equality assertion is global and imposes no proximity to an
ellipsoid.

\subsection*{Previous results and methods}

One approach studies a body together with its second projection body.
The class-reduction inequality $R_n(\Pi K)\le R_n(K)$, together with
its equality condition, reduces the study of minimizers to bodies
homothetic to $\Pi^2K$; see \cite[Sections~1--2]{Saroglou2015}.
Saroglou and Zvavitch proved a local minimum theorem under the
hypothesis that the density of the surface-area measure, after an
invertible linear change of variables, is sufficiently close to that
of the ball in $L^\infty$ \cite[Theorem~1.3]{SaroglouZvavitch2017}.
Ivaki established local rigidity for $C^2$ solutions of
$\Pi^2K=cK$ near the ball after an invertible linear change of
variables \cite{Ivaki2018}. These results use the harmonic structure
of the projection operator and the special role of the directions
arising from linear changes of coordinates.

Chen, Feng, Li, Xi, and Xu proved the unrestricted three-dimensional
inequality, including ellipsoid equality
\cite[Theorem~1.1]{ChenEtAl2026}. Mielke-Sulz proved the conjecture
for bodies of revolution in every dimension $n\ge3$
\cite{MielkeSulz2026}. Theorem~\ref{thm:main} treats arbitrary convex
bodies in the remaining dimensions. Combining it with the separately
proved three-dimensional theorem of Chen et al. establishes the
conjecture for every $n\ge3$; that external theorem is not used in our
proof for $n\ge4$.

Lutwak reformulated the conjecture as a lower bound for a two-body
integral with kernel $|u\cdot v|$ against surface-area measures
\cite{LutwakPetty1990}; see the explicit formulation in
\cite[Introduction, item~(i)]{Saroglou2015}.
Our reduction uses the same pairing for a broader class of positive
measures satisfying a lower bound on integrals of support functions.

Our proof uses the classical first-variation and spherical-transform
frameworks, but requires estimates and an affine normalization that
apply to arbitrary bodies. The Wulff variations are closely related
to the Aleksandrov variations for dual curvature measures developed
by Huang, Lutwak, Yang, and Zhang
\cite{HuangLutwakYangZhang2016}. The variational representation of a
support functional belongs to the $L_p$ dual Minkowski framework;
compare Huang and Zhao \cite{HuangZhao2018}. We prove the needed
representation directly, including its logarithmic endpoint.
The cosine and Funk spectra are classical
\cite{OlafssonPasqualeRubin2013}; we derive their multipliers in our
normalization. The further tasks are a strict estimate for all norms
and a continuous choice of variational minimizers through changes of
matrix rank. The latter permits a global choice of affine coordinates,
where proximity to the ball is unavailable.

\subsection*{The argument}

The proof passes from a body to a positive measure on vectors.
After translating and normalizing $|K|=\kappa_n$, the first variation
of volume gives a bounded, spanning probability measure $\eta_K$ on
$\R^n$ such that
\[
 \int h_M\,d\eta_K\ge (|M|/\kappa_n)^{1/n}
\]
for every origin-symmetric convex body $M$.
The same measure represents $h_{\Pi K}(y)$ as a fixed multiple of
$\int|x\cdot y|\,d\eta_K(x)$.
Theorem~\ref{thm:bilinear} bounds the pairing
\[
 \iint|x\cdot y|\,d\eta(x)d\zeta(y)
\]
for any two bounded, spanning positive measures satisfying this
support-functional inequality. Applying the result to the measures
of $K$ and its volume-normalized projection body yields the desired
lower bound. Sharpness in the pairing returns equality to the
homothety case of the first variation of volume.

Two ingredients have statements independent of projection bodies.
Theorem~\ref{thm:norm-estimate} bounds the higher spherical harmonics
of a power of an arbitrary norm by a difference of its spherical
means. The estimate is strict except for the Euclidean norm, and no
smoothness is needed. Proposition~\ref{prop:representation} represents
a positive support functional using the gradient of a variationally
selected norm. A spherical sign test bounds the pairing below by a
bilinear form involving the cosine and Funk transforms. The strict norm estimate
controls the contribution of harmonics of degree at least four.

The degree-two component must be removed by choosing affine
coordinates; it is not a zero eigenspace of the sign-test operator.
Proposition~\ref{prop:position} makes this choice by extending the
family of variational minimizers continuously to singular matrices.
At that boundary the gauges become seminorms. An explicit calculation
in their missing directions gives an inward-pointing matrix field,
and a fixed-point argument yields an invertible matrix at which the
degree-two component vanishes. This position is needed for only one
of the two measures. The other is transformed contragrediently, which
preserves their pairing. These constructions and the strict norm
estimate also supply the equality case.

Combining Theorem~\ref{thm:main} with the separate three-dimensional
result of Chen et al. gives several sharp inequalities in dimensions
$n\ge3$. Lutwak's implication yields Euclidean-ball lower bounds for
lower-degree projection-body ratios, in which intrinsic volumes replace
the volumes of shadows. Haddad's $p=1$ implication bounds the integral of
the absolute determinant of gradients below by the product of critical
Lebesgue norms. Its
diagonal case strengthens
the Sobolev--Zhang affine $L^1$ inequality. The theorem also gives the
sharp Holmes--Thompson isoperimetric inequality in normed spaces, where
equality requires an ellipsoidal norm ball and a body homothetic to it.
The full statements below distinguish the hypotheses and equality
conclusions of these consequences.

Section~\ref{sec:gauges} establishes the volume and gauge facts used in
the reduction. Section~\ref{sec:spherical} derives the sign test and
its harmonic multipliers. Section~\ref{sec:norm} proves the strict norm
estimate, including the scalar inequalities in every dimension.
Section~\ref{sec:position} constructs the variational representation
and affine position. Section~\ref{sec:conclusion} proves the bilinear
inequality and deduces Theorem~\ref{thm:main}.
Section~\ref{sec:consequences} develops its consequences. We use the classical
Brunn--Minkowski inequality with its homothety equality case and
Brouwer's fixed-point theorem in their stated forms; the first-variation
reduction, harmonic estimates, variational construction, and equality
argument are proved below.

\subsection*{Notation}

Throughout Sections~\ref{sec:gauges}--\ref{sec:conclusion}, $n\ge4$,
$p=n-1$, and $\sigma$ is uniform probability on
$S^{n-1}$. The symbol $\E$ denotes integration against $\sigma$, and
\[
 b_n=\E|u_1|.
\]
Repeated spherical integrals use independent directions.
A symmetric body always means an origin-symmetric body.
When a function on the sphere is used as a degree-one homogeneous test
on vectors, it is extended by positive homogeneity and given value zero
at the origin. Euclidean gradients are denoted by $\nabla$; the
tangential gradient and Laplacian are $\gradS$ and $\lapS$.
